% =============================================================================
\section{Binding code}
\label{sec:code}
% =============================================================================
Boost.Python allows users to expose C++ classes and functions to
Python using nothing more than a C++ compiler. It exploits
metaprogramming techniques to implement a rich set of features and
high-level user-friendly interface. Boost.Python extracts as much as
information as possible from the source code to be wrapped. This
approach is referred to as \emph{user guided wrapping}. Additional
information that cannot be automatically deduced is explicitly
supplied by the user. The interface specification is written in the
same full-featured C++ language as the code being exposed, which has
the potential of generating efficient binding. In this section we
describe some of the techniques we use to ensure that.

% -----------------------------------------------------------------------------
\subsection{Return value policy}
\label{sec:code:rvp}
% -----------------------------------------------------------------------------
Python and C++ use fundamentally different ways of managing the memory
and lifetime of objects. This can lead to issues when creating
bindings for functions that return a non-trivial type. Just by looking
at the type, it is unclear whether Python should take charge of the
returned value and eventually free its resources, or whether this is
handled by the C++ code. When writing code in C++, it is usually
considered a good practice to use smart pointers, which exactly
describe ownership semantics. Still, even good C++ interfaces use raw
references and pointers sometimes. In some cases, in order to assure
proper memory management of a return value from a function, explicitly
specifying a return value policy is needed. The policy may depend on
whether the returned object is a newly created object, a reference to
some already existing (internal) object, or a const reference to some
existing object. Additionally, in some cases we need to tie the
lifetime of the result to the lifetime of the arguments, or the other
way around.

For some functions of kernel objects the kind of return value depends
on the kernel type; thus, the return value policy needs to be set
accordingly. For example, when
\cppLstinline{Exact\_predicates\_inexact\_constructions\_kernel} is
the kernel type, the return type of some functions that return a
coordinate or a coefficient is a const reference, while when the
kernel type is
\cppLstinline{Exact\_predicates\_exact\_constructions\_kernel} the
same functions return an object by value. This is why for these
functions a \myLstinline{Kernel_return_value_policy} type is being
passed as the return value policy. The type is defined differently
based on the kernel being exposed.

%% All functions and classes related to some package were exposed inside
%% a \myLstinline{export_[package]()} function. All such calls are
%% happening in \myLstinline{export_module.cpp}. We wrap each call with a
%% local (C++) scope and call the \myLstinline{SET_SCOPE(module_name)}
%% macro, which causes every exposed function/class inside this local
%% scope to be exposed inside the \myLstinline{module_name} submodule of
%% the CGALPY module. This prevents name clashes of classes with the same
%% name in different packages.

%% % -----------------------------------------------------------------------------
%% \subsection{Avoiding Code Duplication}
%% \label{sec:code:templates}
%% % -----------------------------------------------------------------------------
%% Some types have the exact same nested types and member functions
%% exposed, e.g., \cppLstinline{Triangulation\_2} and
%% \cppLstinline{Delaunay\_triangulation\_2} class templates each have
%% the method \cppLstinline{dimension()} exposed as the Python attribute
%% \pyLstinline{dimension}. In these cases we use a function template to
%% avoid code duplication. Following the example above, we have the
%% function \cppLstinline{export\_triangulation\_2<Triangulation>()},
%% where the template parameter \cppLstinline{Triangulation} is
%% substituted with one of the two-dimensional triangulation types
%% supported; see Table~\ref{tab:tri2:names}. This function template
%% accepts as its single parameter the binding class object, which is
%% constructed separately.

%% \begin{lstlisting}[style=cppNumberedStyle]
%% typedef CGAL::Triangulation_2<Kernel>              Triangulation_2;
%% typedef CGAL::Delaunay_triangulation_2<Kernel>     Delaunay_triangulation_2;

%% auto c0 = bp::class_<Triangulation_2>("Triangulation_2");
%% auto c1 = bp::class_<Delaunay_triangulation_2,
%%                      bp::bases<Triangulation_2>>("Delaunay_triangulation_2");

%% export_triangulation_2<Triangulation_2>(c0);
%% export_triangulation_2<Delaunay_triangulation_2>(c1);
%% \end{lstlisting}

% -----------------------------------------------------------------------------
\subsection{Intersection Detection}
\label{sec:code:intersectionDetection}
% -----------------------------------------------------------------------------
Substitution failure is not an error (SFINAE) refers to a situation in
C++ where an invalid substitution of template parameters is not in
itself an error. In particular, it applies during a process called
\emph{overload resolution}.  In order to compile a function call, the
compiler creates a set of candidate functions the names of which match
the call and that can be accessed by the caller. Then, it reduces the
candidate set to a set of viable functions that includes all function
instances the parameters of which match the call arguments. Finally,
the compiler selects the best match among the viable functions
according to the C++ standard;\footnote{For the full specification of
  the C++ standard see \url{https://isocpp.org/std/the-standard}.} a
flowchart that describes the process is shown in
Figure~\ref{fig:ors}. When the substitution of an explicitly specified
or deduced type for the template parameter fails, the specialization
is discarded from the overload set instead of causing a compilation
error.\footnote{For more information on SFINAE see, e.g.,
  \url{https://en.cppreference.com/w/cpp/language/sfinae}.}

\begin{figure}[!htp]
    \centerline{\begin{tikzpicture}[%
    >=stealth,align=center,every node/.style={font=\sffamily}]
    \tikzstyle{state}=[text centered,node distance=0.5cm,%
      rectangle,rounded corners,fill=white,draw=black,%
      minimum width=1.0cm,minimum height=0.45cm]
    % All state styles
    \tikzstyle{lookup}=[state,bottom color=purple!40]
    \tikzstyle{deduction}=[state,fill=orange!40]
    \tikzstyle{reduction}=[state,fill=yellow!40]
    \tikzstyle{selection}=[state,fill=green!40]
    % All states
    \node (lookup)[lookup]{Name Lookup};
    \node (lookupFailure)[below left=0.5 and -1 of lookup]
      {\textcolor{red}{Error: Lookup failure}};
    \node (deduction)[deduction,below right=1 and -1.2 of lookup]
      {Template Argument Deduction};
    \node (deductionFailure)[below left=0.5 and -2.8 of deduction]
      {\textcolor{red}{Error: Deduction failure}};
    \node (reduction)[reduction,below right=1 and -1.2 of deduction] {Reduction};
    \node (reductionFailure)[below left=0.5 and -1.6 of reduction]
      {\textcolor{red}{Error: Substitution failure}};
    \node (selection)[selection,below right=1 and 0.0 of reduction]
      {Final Selection};
    \node (success)[below right=0.5 and -0.7 of selection]{\textcolor{blue}{Best Match}};
    \node (slectionFailure)[below left=0.5 and -1.6 of selection]
      {\textcolor{red}{Error: Ambiguity failure}};
    % Edges
    \draw [thick,->,>=Stealth] (lookup) -- (deduction);
    \draw [thick,->,>=Stealth] (lookup) -- (lookupFailure);
    \draw [thick,->,>=Stealth] (deduction) -- node[anchor=west]{Non-empty candidate Set}(reduction);
    \draw [thick,->,>=Stealth] (deduction) -- (deductionFailure);
    \draw [thick,->,>=Stealth] (reduction) -- node[anchor=west]{Non-empty viable Set} (selection);
    \draw [thick,->,>=Stealth] (reduction) -- (reductionFailure);
    \draw [thick,->,>=Stealth] (selection) -- (success);
    \draw [thick,->,>=Stealth] (selection) -- (slectionFailure);
  \end{tikzpicture}}
  \caption{Overload resolution steps}
  \label{fig:ors}
\end{figure}

The function \cppLstinline{CGAL::do_intersect(o1, o2)}, provided by
the \iiDandiiiDGeometryKernelPackage{} package (see
Section~\ref{ssec:modules:kernel}), is overloaded with several
implementations that handle different combination of types of
arguments; however, not every combination of two types is implemented
in this package. For example, the intersection of a line segment and a
circle in the general case is a point with algebraic coordinates, and
thus not supported by the package.\footnote{For a complete list of
  valid combinations of \cppLstinline{do_intersect()} argument types
  refer to the reference manual at
  \url{https://doc.cgal.org/latest/Kernel_23/group__do__intersect__linear__grp.html}.}
We use SFINAE to generate bindings for all supported combinations
while avoiding getting compilation errors for unsupported combinations
as follows. Consider two types \cppLstinline{Type1} and
\cppLstinline{Type2}, and assume that the
\cppLstinline{CGAL::do_intersect(Type1 o1, Type2 o2)} overloaded
version is not implemented but all other three are, and refer to
Listing~\ref{lst:do_intersect}.

\begin{lstfloat}
  \begin{lstlisting}[style=cppFramedStyle,basicstyle=\normalsize,%
                     label={lst:do_intersect},%
                     caption={Automatic generation of binding  for \cppLstinline{CGAL::do_intersect()} exploiting SFINAE}]
template <typename, typename> void bind_do_intersect_2T(...) {}

template <typename T1, typename T2>
void bind_do_intersect_2T(decltype(CGAL::do_intersect<Kernel>(T1(), T2()))) {
  bp.def<bool(const T1&, const T2&)>("do_intersect", &CGAL::do_intersect<Kernel>);
}

template <typename T1>
void bind_do_intersect_1T() {
  bind_do_intersect_2T<T1, Type1>(true);
  bind_do_intersect_2T<T1, Type2>(true);
}

void bind_do_intersect() {
  bind_do_intersect_1T<Type1>();
  bind_do_intersect_1T<Type2>();
}
  \end{lstlisting}
\end{lstfloat}

The function \cppLstinline{bind_do_intersect()} calls
\cppLstinline{bind_do_intersect_1T<T1>()} twice; in the first call
\cppLstinline{T1} is substituted with \cppLstinline{Type1} (Line~15)
and in the second call \cppLstinline{T1} it is substituted with
\cppLstinline{Type2} (Line~16).  The function template
\cppLstinline{bind_do_intersect_1T()} calls a third function, called
\cppLstinline{bind_do_intersect_2T()}, twice; first it calls
\cppLstinline{bind_do_intersect_2T<T1, Type1>(true)} (Line~10); then,
it calls \cppLstinline{bind_do_intersect_2T<T1, Type2>(true)}
(Line~11). The function \cppLstinline{bind_do_intersect_2T()} is
overloaded with two implementations. As a consequence, assuming that
the code compiles, one overload of
\cppLstinline{bind_do_intersect_2T()} is selected for every
permutation of (\cppLstinline{Type1}, \cppLstinline{Type2}). The
primary implementation of \cppLstinline{bind_do_intersect_2T()}
(Line~1) uses \emph{variadic arguments} and serves as a fall-through
when the evaluation of the second implementation fails during the
resolution process.\footnote{For more information on \emph{variadic
    arguments} see
  \url{https://en.cppreference.com/w/cpp/language/variadic_arguments}.}
The parameter type of the second implementation is defined as the
return type of \cppLstinline{CGAL::do_intersect<Kernel>(T1(),
  T2())}. When \cppLstinline{T1} and \cppLstinline{T2} are substituted
with \cppLstinline{Type1} and \cppLstinline{Type2}, respectively, the
evaluation fails, and the blank implementation is selected, as this is
the only candidate left. In all other three cases the evaluation
succeeds and the second implementation is selected, as functions that
accept variadic parameters have a lower rank for the purpose of
overload resolution.
% The return type of \cppLstinline{CGAL::do_intersect<Kernel>(T1(),
% T2())} is resolved at compile time.
Observe that in both calls to \cppLstinline{bind_do_intersect_2T()}
(Lines~10-11) we pass a \cppLstinline{bool} argument of value
\cppLstinline{true}. This value can be \cppLstinline{false} just as
well, but its type, that is \cppLstinline{bool}, must match the return
type of \cppLstinline{CGAL::do_intersect<Kernel>(T1(), T2())} when
defined. The result is not only compact code that is easy to maintain,
but also code that supports bindings for every potential
\cppLstinline{do_intersect()} overload that might be implemented in
the future.

% -----------------------------------------------------------------------------
\subsection{Minkowski Sum Construction}
\label{sec:code:minkSumComputation}
% -----------------------------------------------------------------------------
The situation presented by the function template
\cppLstinline{CGAL::minkowski_sum_2()} that applies the
\emph{convex-decomposition approach} (see
Section~\ref{ssec:modules:ms2}) is a bit more complicated. The
function template shown in Line~1 of
Listing~\ref{lst:minkSumConvexDecomposition} has 10 instances of valid
combinations of argument types. There is another set, of the same
cardinality, of function templates that also accept a forth traits
parameter. The number of instances of the function template shown in
Line~8 is 36, and also here there is another set, of the same
cardinality, of function templates that accept an additional traits
parameter. Consider the function template shown in Line~1. We need to
generate bindings for instances of this function only for valid
combinations of \cppLstinline{PolygonType1},
\cppLstinline{PolygonType2}, and
\cppLstinline{PolygonConvexDecomposition_2} types. Here after and in
the code listing we use the respective short names \cppLstinline{PT1},
\cppLstinline{PT2}, and \cppLstinline{PP} instead. A valid combination
exists only if \cppLstinline{PP} is a type of a polygon-partition
unary function that can be applied to a polygon of type
\cppLstinline{PT1} and to a polygon of type \cppLstinline{PT2}. Some
supported polygon-partition functions cannot be applied to polygon
with holes, hence invalid combinations. In particular, if
\cppLstinline{pp} is a polygon-partition function, the calls
\cppLstinline{pp(pgn1, it)} and \cppLstinline{pp(pgn2, it)} must be
valid, where \cppLstinline{pgn1} and \cppLstinline{pgn2} are input
polygons of types \cppLstinline{PT1} and \cppLstinline{PT2},
respectively, and \cppLstinline{it} is an output iterator of a
container of the resulting polygons. Dereferencing the iterator must
yield a type that can represent a simple polygon. If \cppLstinline{PT}
is identical to \cppLstinline{Polygon_2}, then
\cppLstinline{std::list<PT>::iterator} can serve as an output
iterator; otherwise, \cppLstinline{PT} must represent a polygon with
holes, and in this case
\cppLstinline{std::list<PT::Polygon_2>::iterator} can serve as an
output iterator. We use SFINAE yet again to make this distinction and
we use SFINAE one more time to generate binding for
\cppLstinline{CGAL::minkowski_sum_2()} only for valid combinations of
argument types.

\begin{lstfloat}
  \begin{lstlisting}[style=cppframedStyle,basicstyle=\normalsize,%
                     label={lst:minkSumConvexDecomposition},%
                     caption={Signatures of function templates that compute the 2D Minkwoski sum using the convex-decomposition approach.}]
template <typename Kernel, typename Container,
          typename PolygonConvexDecomposition_2>
Polygon_with_holes_2<Kernel, Container>
minkowski_sum_2(const PolygonType1<Kernel, Container>& p,
                const PolygonType2<Kernel, Container>& q,
                const PolygonConvexDecomposition_2& decomp)

template <typename Kernel, typename Container,
          typename PolygonConvexDecompositionP_2,
          typename PolygonConvexDecompositionQ_2>
Polygon_with_holes_2<Kernel, Container>
minkowski_sum_2(const PolygonType1<Kernel, Container>& p,
                const PolygonType2<Kernel, Container>& q,
                const PolygonConvexDecompositionP_2& decomp_p,
                const PolygonConvexDecompositionQ_2& decomp_q)
  \end{lstlisting}
\end{lstfloat}

First, we introduce a class template called \cppLstinline{target}; see
Line~3 in Listing~\ref{minkSumResolution}. it is parameterized with a
named parameter \cppLstinline{T} and an unnamed parameter that
defaults to \cppLstinline{void}. It delegates the type
\cppLstinline{std::list<T>::iterator}. We also introduce a
specialization (Line 7) that delegates the type
\cppLstinline{std::list<T::Polygon_2>::iterator}. The second parameter
is substituted with a type provided by a utility class template called
\cppLstinline{has<>} (Line~1). If the type \cppLstinline{T::Polygon_2}
is undefined, the evaluation of \cppLstinline{has<T::Polygon_2>}
fails, and in turn the instance of the \cppLstinline{Target<>}
specialization is discarded from the overloaded resolution
set. Otherwise, \cppLstinline{has<T::Polygon_2>} successfully
evaluates (to \cppLstinline{void}). In this case the
\cppLstinline{Target<>} specialization remains a viable option, and it
is selected over the primary implementation, since specializations are
ranked higher. Observe that the type delegated by the
\cppLstinline{has<>} class template and the type of the unnamed
template parameter of the primary implementation of
\cppLstinline{target<>} must match for the ordering to apply. These
types are arbitrarily chosen to be \cppLstinline{void}.
%
We need to capture valid combinations of three types. Recall that the
code in Listing~\ref{lst:do_intersect} captures valid combinations of
two types. Following the idea presented in the previous section, we
introduce three template functions, where
\cppLstinline{bind\_mink\_sum\_3T()} is the last one called; see
Line~12 and Line~15 in Listing~\ref{minkSumResolution} for the primary
and specialized implementation, respectively. The polygon-partition
function returns an output iterator that points to next to the last
element of the container. Passing an argument, the type of which
matches the type of the returned value of the polygon-partition
function, to the function \cppLstinline{bind\_mink\_sum\_3T()} is
complicated, as this type is not a simple type (e.g.,
\cppLstinline{bool}) any more. Instead, we introduce two unnamed
template parameters with default values, that evaluate to the return
type of the polygon-partition function when applied to a polygon of
type \cppLstinline{PT1} or type \cppLstinline{PT2}, respectively; see
Lines~16 and 17. Observe that the outcome of the evaluation is
irrelevant. The only thing that matters is weather the evaluation
succeeds---it succeeds only if the matching function is defined.

\begin{lstfloat}
  \begin{lstlisting}[style=cppFramedStyle,basicstyle=\normalsize,
                     label=minkSumResolution,
                     caption={Automatic generation of binding for \cppLstinline{CGAL::minkowski_sum_2} exploiting SFINAE.}]
template <typename T> class has { typedef void type; };

template <typename T, typename = void> struct target {
  typedef typename std::list<T>::iterator type;
};

template <typename T>
struct target<T, typename has<typename T::Polygon_2>::type> {
  typedef typename std::list<typename T::Polygon_2>::iterator type;
};

template <typename T1, typename T2, typename T3>
void bind_mink_sum_3T(...) {}

template <typename PT1, typename PT2, typename PP,
          typename = decltype(PP()(PT1(), typename target<PT1>::type())),
          typename = decltype(PP()(PT2(), typename target<PT2>::type()))>
void bind_mink_sum_3T(bool) {
  bp::def<Polygon_with_holes_2(const PT1&, const PT2&, const PP&)>
    ("minkowski_sum_2", &CGAL::minkowski_sum_2<Kernel, Point_2_container, PP>);
}
  \end{lstlisting}
\end{lstfloat}

%% templated with some supported type
%% {T1}. \myLstinline{bind_do_intersect_1T()}, templated with
%% \myLstinline{T1}, calls \myLstinline{bind_do_intersect_2T} templated
%% with \myLstinline{T1} and some supported type \myLstinline{T2}, and a
%% \myLstinline{bool} argument.  The declaration of
%% \myLstinline{bind_do_intersect_2T}, templated with both
%% \myLstinline{T1} and \myLstinline{T2} is only valid when templated
%% with a combination of types for which an overloaded implementation of
%% \myLstinline{CGAL::do_intersect} exists. Otherwise, due to SFINAE, the
%% blank implementation template of \myLstinline{bind_do_intersect_2T}
%% expecting (...) as an argument will be invoked (it would normally not
%% be invoked as variadic parameters have the lowest rank for the purpose
%% of overload resolution).

\noindent%
We exploit SFINAE to generate bindings of other types as described in the following sections.

% -----------------------------------------------------------------------------
\subsection{Arrangement Extension}
\label{sec:code:ae}
% -----------------------------------------------------------------------------
In many applications it is necessary to extend the type that
represents certain \cgal{} data structures with some types that
represent new properties. In \cgal{} it is more convenient, and thus
more common, extending the types that represent the topological
features of the data structures (rather then the types that represent
the geometric features), e.g., the DCEL of the arrangement data
structure; see Section~\ref{ssec:modules:aos2}. When developing code
in pure C++, there is no restriction on the property types used to
extend the data structure. However, when the data structure interface
is wrapped with Python bindings, the property type must be known when
the bindings are compiled. Nevertheless, we support extensions with
generic Python objects. If a certain feature must be extended with a
certain property, the type of which is defined in the C++ code, it is
always possible to use an external property map, where the properties
are indexed by a Python object, such as an integer or a string. When
generating the bindings the user must set the flag that enables the
extension; see Table~\ref{tab:aos}. Attaching an index to and
extracting an index from a feature is done via the
\cppLstinline{set_data()} and \cppLstinline{data()} exposed
functions. Each cell type, namely, \cppLstinline{Vertex},
\cppLstinline{Halfedge}, or \cppLstinline{Face}, can be extended
independently, which implies that there are eight plausible
combinations that the user can select from to form the final DCEL type
of the arrangement.

Extending the arrangement features is useful for many applications,
but essential for applications that compute the overlay of
arrangements as discussed in the next Section.

% Add a discussion about extanding the triangulation data structure in 2D and 3D
% Change the title to be more general to reflect other extensions

% -----------------------------------------------------------------------------
\subsection{Arrangement Overlay Traits}
\label{sec:code:overlay}
% -----------------------------------------------------------------------------
The map overlay of two arrangements $\calA_1$ and $\calA_2$,
conveniently referred to as the \textcolor{red}{red} and
\textcolor{blue}{blue} arrangements, is a third arrangement $\calA$,
such that there is a cell $c$ in $\calA$ iff there are cells $c_1$ and
$c_2$ in $\calA_1$ and $\calA_2$, respectively, such that $c$ is a
maximal connected component of $c_1 \cap c_2$.
% Arrangements are subdivisions of some ambient space (see Section~\ref{ssec:modules:aos2})
Computing the overlay of two arrangements is useful for many
applications, including Boolean operations; see
Section~\ref{ssec:modules:bso2}. Indeed, the \iiDArrangementsPackage{}
package provides function templates that computes the overlay of two
arrangements, namely overloaded \cppLstinline{CGAL::overlay()}. In
particular, the call \cppLstinline{CGAL::overlay(arr_r, arr_b, arr)}
computes the overlay of the arrangements \cppLstinline{arr_r} and
\cppLstinline{arr_b} and stores the result in the arrangement
\cppLstinline{arr}. All three types must be instances of the
\cppLstinline{Arrangement_2<Traits, Dcel>} class template; see
Section~\ref{ssec:modules:aos2}. When the function is used in pure C++
code, the geometry-traits classes that substitute the
\cppLstinline{Traits} template parameters of the input arrangements
must be convertible to the geometry-traits class of the resulting
arrangement. We only wrap instances of the overlay function with
Python bindings, where all three arrangements use the exact same
geometry-traits classes and the same DCEL structures, and thus the
types of the three arrangements of the exposed
\cppLstinline{CGAL::overlay()} function are identical.

The \cppLstinline{CGAL::overlay()} function template is overloaded
with a variant that accepts four arguments. The last argument,
referred to as the overlay traits, enables the update of cells of the
output arrangement, and in particular their properties. Using the
overlay traits is typically accompanied with the extension of one or
more types of the arrangement cells that hold the properties. The
overlay traits must model the \cppLstinline{OverlayTraits} concept,
which requires the provision of ten functions that handle all possible
overlapping cases as listed below. Let $v_r$, $e_r$, and $f_r$ denote
input \textcolor{red}{red} cells, i.e.,~a vertex, an edge, and a face,
respectively, $v_b$, $e_b$, and $f_b$ denote input
\textcolor{blue}{blue} cells, and $v$, $e$, and $f$ denote output
cells.
\begin{enumerate}
\item A new vertex $v$ is induced by coinciding vertices $v_r$ and $v_b$.
\item A new vertex $v$ is induced by a vertex $v_r$ that lies on an edge $e_b$.
\item An analogous case of a vertex $v_b$ that lies on an edge $e_r$.
\item A new vertex $v$ is induced by a vertex $v_r$ that is contained in a face $f_b$.
\item An analogous case of a vertex $v_b$ contained in a face $f_r$.
\item A new vertex $v$ is induced by the intersection of two edges $e_r$ and $e_b$.
\item A new edge $e$ is induced by the (possibly partial) overlap of two edges $e_r$ and $e_b$.
\item A new edge $e$ is induced by the an edge $e_r$ that is contained in a face $f_b$.
\item An analogous case of an edge $e_b$ contained in a face $f_r$.
\item A new face $f$ is induced by the overlap of two faces $f_r$ and $f_b$.
\end{enumerate}

Evidently a custom C++ overlay traits cannot be defined in Python;
more precisely, the traits type must be known when the bindings are
compiled. We introduce and expose with Python bindings two models of
the \cppLstinline{OverlayTraits} concept. The most general among the
two, called \cppLstinline{Arr_overlay_traits}, defines ten functions
that correspond to the list above, each accepting three arguments,
namely, two objects that represent input cells, and one object that
represents an output cell. By default these functions are idle. A user
(that is, a Python programmer) can override any subset of these
functions. One constructor of this model accepts all the ten functions
at once. Another constructor accepts a single function, which
corresponds to Function (10) in the list above.

In the following we assume that the face type of the arrangement
% \cppLstinline{Arrangement_2} instance
is extended. The code excerpt in Listing~\ref{lst:twoFaceExtend} constructs two arrangements with two faces each. For each arrangement the properties of the unbounded face and the bounded face are initialized with Python integer objects~$0$ and~$1$, respectively.
%
The code excerpt in Listing~\ref{lst:overlayTraits} computes the overlay of the two arrangements constructed by the code in Listing~\ref{lst:twoFaceExtend} using the overlay traits. The property of each face of the resulting arrangement is updated as part of the overlay computation to indicate the number of overlapping bounded faces; see Figure~\ref{fig:faceExt}.
%
Observe that while the computation of the overlay is carried out by the compiled C++ code, the summation is carried out by a Python \pyLstinline{lambda} function, which accepts as the third argument the face to be updated. By default, arguments are copied into new Python objects. We pass the new face by reference, overriding the default, to enforce persistent updates.

\begin{lstfloat}
  \lstinputlisting[style=pythonFramedStyle,basicstyle=\normalsize,
                   firstline=14,lastline=39,
                   label={lst:twoFaceExtend},
                   caption={Constructing two arrangements with their faces extended and initializing their data.}]
                  {../../../src/python_scripts/cgalpy_examples/overlay.py}
\end{lstfloat}

\begin{lstfloat}
  \begin{lstlisting}[style=pythonFramedStyle,basicstyle=\normalsize,
                     label={lst:overlayTraits},
                     caption={Computing the map overlay of two arrangements using the general overlay traits.}]
traits = Aos2.Arr_overlay_traits(lambda f1,f2,f: f.set_data(f1.data()+f2.data()))
Aos2.overlay(arr1, arr2, result, traits)
  \end{lstlisting}
\end{lstfloat}

\begin{figure}[!htp]
  \captionsetup[subfigure]{justification=centering}
  \centering%
  \subfloat[]{\label{fig:faceExt}
    \begin{tikzpicture}
      \fill[polyALightColor] (0,0) rectangle (2,2);
      \fill[polyBLightColor,opacity=0.5] (1,1) rectangle (3,3);
      \draw[thick,polyADarkColor] (0,0) rectangle (2,2);
      \draw[thick,polyBDarkColor] (1,1) rectangle (3,3);
      \node at (2.5,0.5) {[0]};
      \node at (0.5,0.5) {[1]};
      \node at (2.5,2.5) {[1]};
      \node at (1.5,1.5) {[2]};
      \node[point] at (0,0) {};
      \node[point] at (2,0) {};
      \node[point] at (0,2) {};
      \node[point] at (2,2) {};
      \node[point] at (1,1) {};
      \node[point] at (3,1) {};
      \node[point] at (1,3) {};
      \node[point] at (3,3) {};
      \node[point] at (2,1) {};
      \node[point] at (1,2) {};
    \end{tikzpicture}}\quad\quad
  \subfloat[]{\label{fig:vertexExt}
    \begin{tikzpicture}[scale=0.2]
      \draw[thick,polyADarkColor] (3,4) circle (5);
      \draw[thick,polyADarkColor] (-3,-4) circle (5);
      \draw[thick,polyBDarkColor] (4,-3) circle (5);
      \draw[thick,polyBDarkColor] (-4,3) circle (5);
      \node[point=red!50!white,label={[label distance=-1pt]180:4}] at (-8, -4) {}; % 4
      \node[point=green,label={[label distance=-1pt]180:6}] at (-7, -1) {};        % 6
      \node[point=red!50!white,label={[label distance=-1pt]180:2}] at (-9, 3) {};  % 2
      \node[point=red!50!white,label={[label distance=-1pt]180:4}] at (-2, 4) {};  % 4
      \node[point=green,label={[label distance=-1pt]90:6}] at (-1, 7) {};          % 6
      \node[point=red!50!white,label={[label distance=-1pt]180:2}] at (-1, -3) {}; % 2
      \node[point=blue,label={[label distance=-1pt]0:12}] at (0, 0) {};            % 12
      \node[point=green,label={[label distance=-1pt]-90:6}] at (1, -7) {};         % 6
      \node[point=red!50!white,label={[label distance=-1pt]0:2}] at (1, 3) {};     % 2
      \node[point=red!50!white,label={[label distance=-1pt]0:4}] at (2, -4) {};    % 4
      \node[point=green,label={[label distance=-1pt]0:6}] at (7, 1) {};            % 6
      \node[point=red!50!white,label={[label distance=-1pt]0:4}] at (8, 4) {};     % 4
      \node[point=red!50!white,label={[label distance=-1pt]0:2}] at (9, -3) {};    % 2
    \end{tikzpicture}}
  \caption[]{%
    (\subref{fig:faceExt}) The arrangement faces are extended with integers that indicate the number of overlapping bounded faces, shown in brackets.
    (\subref{fig:vertexExt}) The arrangement vertices are extended with integers that indicate the weighted degree of the vertices.}
  \label{fig:extend}
\end{figure}

In many cases, such as in the example above, the property of the new
cell depends solely on the properties of the overlapping two cells. To
this end we introduce (and expose) a second model of the
\cppLstinline{OverlayTraits} concept called
\cppLstinline{Arr_overlay_function_traits}. This model also defines
ten functions that correspond to the list above. However, here each
function accepts the two Python objects that extend the overlapping
cells as input arguments and returns the Python object that extends
the resulting cell. The code excerpt in
Listing~\ref{lst:overlayFunctionTraits} has the same effect the code
in Listing~\ref{lst:overlayTraits} has, but is more compact.

\begin{lstfloat}
  \begin{lstlisting}[style=pythonFramedStyle,basicstyle=\normalsize,
                     label={lst:overlayFunctionTraits},
                     caption={Computing the map overlay of two arrangements using the functional overlay traits.}]
traits = Aos2.Arr_overlay_function_traits(lambda x, y: x+y)
os2.overlay(arr1, arr2, result, traits)
  \end{lstlisting}
\end{lstfloat}

Each one of the ten functions supported by the
\cppLstinline{Arr_overlay_function_traits} model returns a Python
object that is passed back to the compiled C++ code. Then, the object
is stored with the new cell. Thus, there is no need to pass arguments
by reference (even thought it could be more efficient in certain
cases); however, the implementation of this model presents a new
coding challenge---if a certain cell is not extended, attempting to
access its non-existing extension would cause a compilation error.
%
Assume that only the vertex type of the \cppLstinline{Arrangement_2}
instance is extended. The Python code excerpt in
Listing~\ref{lst::overlayCircleSegment} computes the overlay of two
arrangements and updates the property of each vertex of the resulting
arrangement to indicate the weighted degree of the vertex, where blue
incident edges weigh twice as much as red incident edges; see
Figure~\ref{fig:vertexExt}. We set all the six functions that update
output vertices in the \cppLstinline{Arr_overlay_function_traits}
object using dedicated setters. The name of a setter matches the
pattern
\lstinline[language={C++},mathescape=true,columns=fixed]{set_$c_rc_b$_$c$()},
where each of $c_r$, $c_b$, and $c$ can be substituted with
\cppLstinline{v}, \cppLstinline{e}, or \cppLstinline{f}. $c_r$ and
$c_b$ determine the input red and blue cell types, respectively, and
$c$ determines the output cell type.

\begin{lstfloat}
  \begin{lstlisting}[style=pythonFramedStyle,basicstyle=\normalsize,
                     label={lst::overlayCircleSegment},
                     caption={Constructing two arrangements with their vertices extended and computing their map overlay using the overlay function traits.}]
Aos2.insert(arr1,[Curve(Circle(Point(3,4),25)),Curve(Circle(Point(-3,-4),25))])
Aos2.insert(arr2, [Curve(Circle(Point(3,4),25)), Curve(Circle(Point(3,4),25))])
for arr in [arr1, arr2]:
  for v in arr.vertices():
    v.set_data(v.degree())
traits = Aos2.Arr_overlay_function_traits()
traits.set_vv_v(lambda x, y: 2*x+y)
traits.set_ve_v(lambda x, y: 2*x+2)
traits.set_vf_v(lambda x, y: 2*x)
traits.set_ev_v(lambda x, y: 4+y)
traits.set_fv_v(lambda x, y: y)
traits.set_ee_v(lambda x, y: 6)
Aos2.overlay(arr1, arr2, result, traits)
  \end{lstlisting}
\end{lstfloat}

Consider a generic function called \cppLstinline{apply()} that accepts
two objects that represent input cells, namely \cppLstinline{r} and
\cppLstinline{b}, one object that represents an output cell, namely
\cppLstinline{o}, and a function called \cppLstinline{f}. The
objective of \cppLstinline{apply()} is to apply \cppLstinline{f} to
the properties of the cells \cppLstinline{r} and \cppLstinline{b} and
store the return value as the property of the cell
\cppLstinline{o}. Recall that a property of a cell \cppLstinline{c} is
obtained and set via the calls \cppLstinline{c.data()} and
\cppLstinline{c.set_data()}, respectively. If the type of the output
cell is not extended, however, \cppLstinline{apply()} should become
idle. Similarly, if the type of an input cell is not extended, the
\pyLstinline{None} Python object (represented by
\cppLstinline{bp::object}) should be passed to \cppLstinline{f}
instead. We use SFINAE yet again, twice, to address the above. First,
we introduce an overload of \cppLstinline{apply()} that is idle and
serves as a fall-through; see Line~9 in Listing~\ref{lst:data}. When a
call to \cppLstinline{apply()} is made, the idle overload is selected
by the overload resolution process if the output cell does not have a
member called \cppLstinline{set_data()}. If the output cell does have
this member the other implementation is selected and calls
\cppLstinline{data(r)} and \cppLstinline{data(b)} to obtain the
properties of the \textcolor{red}{red} and \textcolor{blue}{blue}
input cells, respectively. Second, we introduce an overload of
\cppLstinline{data()} that serves as a fall-through (Line~2). It does
nothing but return the \pyLstinline{None} Python object. It is
selected if the corresponding cell does not have a member called
\cppLstinline{data()}. If the cell does have this member the other
implementation is selected; it returns the property of the cell.

\begin{lstfloat}
  \begin{lstlisting}[style=cppFramedStyle,basicstyle=\normalsize,
                     label={lst:data},
                     caption={Automatic generation of bindings for the overlay function traits exploiting SFINAE.}]
// Fall-through; A::data() does not exist
template <typename A> bp::object data(...) { return bp::object(); }

// A::data() exists
template <typename A, typename = decltype(std::declval<A>().data())>
const bp::object& data(const A* a) { return a->data(); }

// Fall-through; O::set_data() does not exist
template <typename R, typename B, typename O, typename F> void apply(...) {}

// O::set_data() does exist
template <typename R, typename B, typename O, typename F,
          typename = decltype(std::declval<O>().set_data(std::declval<typename O::Data>()))>
void apply(const R* r, const B* b, O* o, F f) {
  o->set_data(f(data<R>(r), data<B>(b)));
}
  \end{lstlisting}
\end{lstfloat}

% -----------------------------------------------------------------------------
\subsection{Miscelaneous}
\label{sec:code:misc}
% -----------------------------------------------------------------------------
In this section we list additional techniques used by the binding code.

% -----------------------------------------------------------------------------
\subsubsection{Handle, iterators, and Circulators}
\label{sec:code:misc:iters}
% -----------------------------------------------------------------------------
The \cppLstinline{Handle} concept, provided by the
\HandlesandCirculatorsPackage{} package~\cite{cgal:dksy-hc-21},
describes types that are akin to pointers to objects. A handle
provides the dereference operator (\cppLstinline{operator*()}) and the
member access operator (\cppLstinline{->()}). Since Python does not
have any kind of pointers, both a handle to an object in \cgal{} and
the object itself are converted to the same Python object. For
example, both types \cppLstinline{Arrangement_2::Vertex} and
\cppLstinline{Arrangement_2::Vertex_handle} are converted to the
Python type \pyLstinline{Vertex}, which is an attribute of the Python
type \pyLstinline{Arrangement_2}.

The \cppLstinline{Iterator} concept and its refinements describe types
that can be used to identify and traverse the elements of container
that contains seqential data.\footnote{For more information about STL
  iterators, see, e.g.,
  \url{https://en.cppreference.com/w/cpp/iterator}.} An iterator
provides the increment or decrement operators. In \cgal{} quite often
an iterator is convertible to a handle, for example,
\cppLstinline{Arrangement_2::Vertex_iterator} is convertible to
\cppLstinline{Arrangement_2::Vertex_handle}. Therefore, also
\cppLstinline{Arrangement_2::Vertex_iterator} is converted to the
python class \pyLstinline{Vertex}. Every converted iterator is
supplied with both magic functions \pyLstinline{__iter__} and
\pyLstinline{__next__}, thus made a \emph{Python
  iterator}.\footnote{For more information about Python iterators,
  See, e.g., \url{https://docs.python.org/3/c-api/iterator.html}.}
The \pyLstinline{__iter__} function simply returns the input object,
and the \pyLstinline{__next__} function returns the next object in the
traversal order of the corresponding iterator.

Similar to the \cppLstinline{Iterator} concept and its refinements,
the \cppLstinline{Circulator} concept and its refinements, also
provided by the \HandlesandCirculatorsPackage{} package, describe
types that can be used to identify and traverse the elements of
container that contains circular data, for example, the halfedges
incident to a face or the halfedges incident to a vertex in an
arrangement.  A circulator object does not have a past-the-end
value. Instead, the range $[a, b)$ of two circulators $a$ and $b$
  denotes either the empty range, or the sequence of all elements in
  the container. Boost.Python supports several features that aid in
  the wrapping if iterators. Circulators are artificially converted to
  iterators to leverage on those features. In addition, while
  traversing a circulator, an stop-iteration exception is thrown when
  a circular refers to an item that has been traversed already.

% -----------------------------------------------------------------------------
\subsection{Function Accespting Collections of Elements}
\label{sec:code:input}
% -----------------------------------------------------------------------------
The Standard Template Library (STL) provides various type-safe
containers for storing collections of related objects. The
\cppLstinline{Container} concept describes types of objects that store
collections of other objects (their elements).\footnote{For more
  information about STL containers, see, e.g.,
  \url{https://en.cppreference.com/w/cpp/container}.} Container models
are implemented as class templates, which enables great flexibility in
the elements types and in the implemented algorithms that operate on
container objects. Traditionally, implementations of algorithms in C++
manipulate iterators pointing into the containers they operate on. To
date it is popular to pass a collection of elements to a function via
two generic input iterators, where the first points at the first
element of a container and the second points past the end of the
container. An input iterator object, the type of
which is a model of the the \cppLstinline{InputIterator} concept,
\footnote{For more information of output iterators, see, e.g.,
  \url{https://en.cppreference.com/w/cpp/iterator/input_iterator}.}
supports reading from the location obtained via the dereference
operator and can be pre- and post incremented.  The Boost.Range
library provides an alternative modern method, which uses
\emph{ranges}, for applying algorithms on collections of objects. In
particular, it utilizes a new family of concepts that refine a basic
concept called \cppLstinline{Range}. The \cppLstinline{Range} concept
is similar to the \cppLstinline{Container} concept; it requires the
provision of iterators for accessing a half-open range of elements and
provides information about the number of elements in the range. The
new method results in code that is more efficient and more expressive
(and thus more comprehensible).\footnote{See
  \url{https://www.boost.org/doc/libs/1_78_0/libs/range/doc/html/index.html}
  for the documentation of Boost.Range.}  Soon, (C++20) ranges will
become part of the standard.\footnote{For more information about STR
  ranges, see, e.g., \url{https://en.cppreference.com/w/cpp/ranges}.}
The \HandlesandCirculatorsPackage{} package of \cgal{} provides a
variant of a \cppLstinline{Range} concept suitable for \cgal{}. Many
functions in \cgal{} accept as input collections of elements, and some
of them already use this concept alreay. The use of ranges is expected
to grow. For every function in \cgal{} that operates on one or more
collections of elements, and regardless of the interface of the
function, that is, whether a pair of generic input iterators, or a
range, is used to represent each collection, we introduce and expose a
wrapper function that accepts a Python list (of type
\cppLstinline{bp::list}) as an argument;\footnote{For more information
  on Python lists, see, e.g.,
  \url{https://docs.python.org/3/tutorial/datastructures.html\#more-on-lists}.}
the wrapper function calls the original function, properly passing the
input collection; see Listing~\ref{lst:insert} for an example of such
a function that uses iterators; the wrapper function is shown in
Listing~\ref{lst:insertWrapper}; an excerpt code in Python that
exploits the above is shown in Listing~\ref{lst:insertPython}.

\begin{lstfloat}
  \begin{lstlisting}[style=cppFramedStyle,basicstyle=\normalsize,%
                     label={lst:insert},%
                     caption={The signature of the free function that inserts a collections of curves into an arrangement.}]
template <typename InputIterator>
CGAL::insert(Arrangement_2& arr, InputIterator first, InputIterator past_the_end)
  \end{lstlisting}
\end{lstfloat}

\begin{lstfloat}
  \begin{lstlisting}[style=cppFramedStyle,basicstyle=\normalsize,%
                     label={lst:insertWrapper},%
                     caption={A function that wraps the function shown in Listing~\ref{lst:insert} and accepts a Python list that contains the input curves.}]
void insert_curves(Arrangement_2& arr, bp::list& lst) {
  if (! lst) return;
  if (bp::extract<X_monotone_curve_2>(lst[0]).check()) {
    auto begin = bp::stl_input_iterator< X_monotone_curve_2 >(lst);
    auto end = bp::stl_input_iterator< X_monotone_curve_2 >();
    CGAL::insert(arr, begin, end);
  }
  else if (bp::extract<Curve_2>(lst[0]).check()) {
    auto begin = bp::stl_input_iterator< Curve_2 >(lst);
    auto end = bp::stl_input_iterator< Curve_2 >();
    CGAL::insert(arr, begin, end);
  }
}
  \end{lstlisting}
\end{lstfloat}

\begin{lstfloat}
  \begin{lstlisting}[style=pythonFramedStyle,basicstyle=\normalsize,%
                     label={lst:insertPython},%
                     caption={A code sample in Python that inserts three curves into an arrangement.}]
arr = Arrangement_2()
c1 = X_monotone_curve_2(Point_2(0, 0), Point_2(1, 0))
c2 = X_monotone_curve_2(Point_2(1, 0), Point_2(0, 1))
c3 = X_monotone_curve_2(Point_2(0, 1), Point_2(1, 1))
Aos2.insert(arr, [c1, c2, c3])
  \end{lstlisting}
\end{lstfloat}

% -----------------------------------------------------------------------------
\subsection{Function Resulting in Collections of Elements}
\label{sec:code:output}
% -----------------------------------------------------------------------------
Functions may compute collections of elements as results. Using output
iterators enables an efficient and flexible method for passing output
collections from functions. An output iterator object, the type of
which is a model of the the \cppLstinline{OutputIterator}
concept,\footnote{For more information of output iterators, see, e.g.,
  \url{https://en.cppreference.com/w/cpp/iterator/output_iterator}.}
supports writing to the location obtained by the dereference operator
and can be pre- and post-incremented. A function that computes a
collection of elements accepts an output iterator as an argument and
popultes the underlying collection. Moderns C++ compilers support the
efficient transfer of resources from one object of a certain type to
another object of the same type, reffered to as \emph{move
  semantics}.\footnote{For more information of move semantics, see,
  e.g.,
  \url{https://en.cppreference.com/w/cpp/language/move_constructor}.}
If the underlying collection supports move semantics, a function that
computes a collection of elements can simply return the collection,
resulting in an elegant yet efficient code. For every function in
\cgal{} that results in a collection of elements, and regardless of
the interface of the function, that is, whether an output iterator is
used or the collection is returned, we introduce and expose a wrapper
function that returns a Python list (of type \cppLstinline{bp::list});
the wrapper function calls the original function, properly populating
a Python list with the output collection, and finally it returns the
list.

The free function template \cppLstinline{CGAL::decompose()} accepts an
arrangement $\mathcal{A}$ and computes a collection of polymorphic
elements via an output iterator. For each vertex $v$ of $\mathcal{A}$
the output collection contains a pair of features---one that directly
lies below $v$ and another that directly lies above $v$. Let $v$ be a
vertex of $\mathcal{A}$. The feature above (respectively below) $v$
may be one of the following:
\begin{compactitem}
\item Another vertex $u$ having the same $x$-coordinate as $v$.
\item An arrangement edge associated with an $x$-monotone curve that
  contains $v$ in its $x$-range.
\item An unbounded face in case $v$ is incident to an unbounded face,
  and there is no curve lying above (respectively below) it.
\item An empty object, in case $v$ is the lower (respectively upper)
  endpoint of a vertical edge in the arrangement.
\end{compactitem}
Listing~\ref{lst:decompose} shows the signature of the
function. Dereferencing the output iterator must yield an object the
type of which is \cppLstinline{Decompose_result} shown in
Listing~\ref{lst:decomposeResult}.

\begin{lstfloat}
  \begin{lstlisting}[style=cppFramedStyle,basicstyle=\normalsize,%
                     label={lst:decompose},%
                     caption={The signature of the free function that decomposes an arrangement into pseudo trapezoids.}]
template <typename OutputIterator>
CGAL::decompose(Arrangement_2& arr, OutputIterator oi);
  \end{lstlisting}
\end{lstfloat}

\begin{lstfloat}
  \begin{lstlisting}[style=cppFramedStyle,basicstyle=\normalsize,%
                     label={lst:decomposeResult},%
                     caption={The type of an output element of the
                       \cppLstinline{CGAL::decompose()} function shown in
                       Listing~\ref{lst:decompose}.}]
typedef std::pair<Arrangement_2::Vertex_const_handle,
                  std::pair<boost::optional<variant>,
                            boost::optional<variant>>>        Decompose_result;
  \end{lstlisting}
\end{lstfloat}

Populating a standard container, say \cppLstinline{container}, the
elements of which are of the \cppLstinline{Decompose_result}
polymorphic type, with the data computed by
\cppLstinline{CGAL::decompose()} function in C++, can be done, for
example, by passing an object of type
\cppLstinline{std::back_insert_iterator<Decompose_result>} obtained by
the call \cppLstinline{std::back_inserter(container)}, as the output
iterator argument to the function. We introduce a similar iterator
called \cppLstinline{apply_iterator} shown in
Listing~\ref{lst:applyIterator} that can be used to efficiently
populate a Python list instead of a container. The constructor of this
function object accepts a generic unary function and stores it for
later use. It applies the operator on every element that needs to be
appended to the outout Python list. Each element of the list is a
Python tuple of two items, the first is a vertex and the second is
another Python tuple of two items, the first is the cell above the
vertex or \pyLstinline{bp::Object} is none exists and the first is the
cell below the vertex or \pyLstinline{bp::Object} is none exists.

\begin{lstfloat}
  \begin{lstlisting}[style=cppFramedStyle,basicstyle=\normalsize,%
                     label={lst:applyIterator},%
                     caption={An output iterator that applies an operation on
                       every output element.}]
template <typename UnaryOperation>
class apply_iterator {
private:
  UnaryOperation m_op;

public:
  apply_iterator(UnaryOperation op) : m_op(op) {}

  // Make sure the assignment operator is available
  apply_iterator& operator=(const apply_iterator& other) {
    m_op = other.m_op;
    return *this;
  }

  template <typename T>
  const T& operator=(const T& t) const { m_op(t); return t; }

  apply_iterator& operator*() { return *this; }
  apply_iterator& operator++() { return *this; }
  apply_iterator operator++(int) { return *this; }
};
  \end{lstlisting}
\end{lstfloat}

The wrapper function of the \cppLstinline{CGAL::decompose()} function
is shown in Listing~\ref{lst:decomposeWrapper}. It utilizes a helper
function called \cppLstinline{decompose_helper()} (not listed) that
appends the output elements to the outout Python list.

\begin{lstfloat}
  \begin{lstlisting}[style=cppFramedStyle,basicstyle=\normalsize,%
                     label={lst:decomposeWrapper},%
                     caption={A function that wraps the function shown in Listing~\ref{lst:decompose} and returns a Python list that represents a vertical decomposition.}]
bp::list decompose(Arrangement_2& arr) {
  typedef std::pair<Arrangement_2::Vertex_const_handle,
                    std::pair<boost::optional<variant>,
                              boost::optional<variant>>>        Decompose_result;
  bp::list lst;
  auto op = [&] (const Decompose_result& res) { decompose_helper(res, lst); };
  CGAL::decompose(arr, apply_iterator<decltype(op)>(op));
  return lst;
}
  \end{lstlisting}
\end{lstfloat}

The Python code example shown in Listing~\ref{lst:py:decompose}
exploits the above and computes the vertical decomposition shown in
Figure~\ref{fig:decompose}.

%% The output:
%% 0 0
%%   unbounded face
%%   unbounded face
%% 2 0
%%   unbounded face
%%    0 0 3 3
%% 3 3
%%    2 0 4 2
%%   unbounded face
%% 4 2
%%   unbounded face
%%   unbounded face
%% 5 3
%%    4 2 6 0
%%   unbounded face
%% 6 0
%%   unbounded face
%%    5 3 8 0
%% 8 0
%%   unbounded face
%%   unbounded face

\begin{lstfloat}
  \lstinputlisting[style=pythonFramedStyle,basicstyle=\normalsize,%
                   label={lst:py:decompose},%
                   firstline=15,
                   caption={A code sample in Python that decomposes an
                     arrangement.}]
  {../../../src/python_scripts/cgalpy_examples/vertical_decomposition.py}
\end{lstfloat}

\begin{figure}[!htp]
  \centerline{%
    \begin{tikzpicture}
      \draw node[point=white] (x) at (4,2) {};
      %
      \draw node[point] (u1) at (0,0) {};
      \draw node[point=cyan] (u2) at (2,2) {};
      \draw node[point] (u3) at (3,3) {};
      \draw node[point=cyan] (u4) at (5,1) {};
      \draw node[point] (u5) at (6,0) {};
      \draw[blue,thick] (u1)--(u2)--(u3)--(x)--(u4)--(u5);
      %
      \draw node[point] (v1) at (2,0) {};
      \draw node[point=cyan] (v2) at (3,1) {};
      \draw node[point] (v3) at (5,3) {};
      \draw node[point=cyan] (v4) at (6,2) {};
      \draw node[point] (v5) at (8,0) {};
      \draw[blue,thick] (v1)--(v2)--(x)--(v3)--(v4)--(v5);
      %
      \draw[dots=10 per 1cm,very thick] (0,-1)--(u1)--(0,4);
      \draw[dots=10 per 1cm,very thick] (2,-1)--(v1)--(u2)--(2,4);
      \draw[dots=10 per 1cm,very thick] (3,-1)--(v2)--(u3)--(3,4);
      \draw[dots=10 per 1cm,very thick] (4,-1)--(x)--(4,4);
      \draw[dots=10 per 1cm,very thick] (5,-1)--(u4)--(v3)--(5,4);
      \draw[dots=10 per 1cm,very thick] (6,-1)--(u5)--(v4)--(6,4);
      \draw[dots=10 per 1cm,very thick] (8,-1)--(v5)--(8,4);
    \end{tikzpicture}}
  \caption{An arrangement of four line segments and its vertical
    decomposition into pseudo trapezoids, as constructed by the Python
    code shown in Listing~\ref{lst:decomposePython}.  The segments of
    the arrangement are drawn in solid blue lines and the segments of
    the vertical decomposition are drawn in dark dotted lines.}
  \label{fig:decompose}
\end{figure}

%% % -----------------------------------------------------------------------------
%% \subsection{\myLstinline{Python_distance}}
%% \label{sec:code:distance}
%% % -----------------------------------------------------------------------------
%% Similarly to overlay traits, a user may want to define a custom
%% distance class to use with search queries. Here too, as defining a new
%% C++ class is not possible from python, we provide a similar solution:
%% A \myLstinline{Python_distance} class is exposed, which is a model of
%% \myLstinline{GeneralDistance}
%% concept. \url{https://doc.cgal.org/5.2/Spatial_searching/classGeneralDistance.html}
%% It expects 5 functions in its constructor, and those are saved as
%% class members and are used for the respective 5 operations [the
%%   optional are not included] of the concept.

%% Making some things more pythonic:

%% (the extended arrangement data type) Those 10 functions are saved as
%% class members and are then called in the respective 10 [overloaded]
%% versions of \myLstinline{create_vertex}, \myLstinline{create_edge},
%% \myLstinline{create_face object}) (the matching is done by the order
%% of the passed functions and the order of the [overloaded] versions in
%% the official concept reference).

%% This way the user can still define the way the result face/edge/vertex
%% data will be calculated.
